\documentclass[10pt,a4paper]{article}
\usepackage[margin=1.55cm]{geometry}
\usepackage{xcolor}
\usepackage{amsmath,amssymb}
\usepackage{graphicx}
\usepackage{enumitem}
\usepackage{fontspec}
\usepackage[CJKmath=true]{xeCJK}
% CJK main font is resolved at COMPILE time on the actual machine, not baked in at
% generation time, so a .tex generated on one OS still renders on another (the older
% Python-side probe hardcoded one OS's font name into the file and broke when moved).
% Walk a cross-platform priority list and use the first font that exists; AutoFakeBold
% + AutoFakeSlant give bold/italic CJK even on faces that ship no bold/italic cut.
\newcommand{\setCJKto}[1]{\setCJKmainfont{#1}[AutoFakeBold=true,AutoFakeSlant=0.2]}
\IfFontExistsTF{Noto Sans CJK SC}{\setCJKto{Noto Sans CJK SC}}{%        % Linux (fonts-noto-cjk); cross-platform if installed
 \IfFontExistsTF{Source Han Sans SC}{\setCJKto{Source Han Sans SC}}{%   % Adobe 思源黑体; cross-platform
  \IfFontExistsTF{PingFang SC}{\setCJKto{PingFang SC}}{%               % macOS 10.11+
   \IfFontExistsTF{Hiragino Sans GB}{\setCJKto{Hiragino Sans GB}}{%    % older macOS
    \IfFontExistsTF{Microsoft YaHei}{\setCJKto{Microsoft YaHei}}{%     % Windows
     \IfFontExistsTF{WenQuanYi Micro Hei}{\setCJKto{WenQuanYi Micro Hei}}{% % older Linux
      \setCJKto{Songti SC}}}}}}}                                       % macOS bottom fallback
\definecolor{cblue}{HTML}{4C72B0}
\setlength{\parskip}{3pt}\setlength{\parindent}{0pt}
\setlist[enumerate]{topsep=2pt,itemsep=2pt,parsep=1pt,partopsep=0pt}
\setlist[itemize]{topsep=2pt,itemsep=2pt,parsep=1pt,partopsep=0pt}
% Let prose-embedded inline math break across lines and absorb small overflows, so simple
% formulas inserted mid-sentence stop running past the right margin.
\tolerance=1200
\emergencystretch=2.5em
\relpenalty=20
\binoppenalty=40
\hfuzz=1pt
\begin{document}

\section*{面向功能几何反事实的边界得分匹配}
\textbf{方法名称：} Boundary-Local Score Matching (BLSM)
\section*{研究动机}
当间隙、姿态或接触几何发生很小变化时，可行动作集合可能改变，而指令、物体和动作的大部分仍保持不变。普通模仿学习不会直接告诉连续动作视觉-语言-动作（VLA）策略这个边界在哪里。全局可行性监督可以拒绝无效轨迹，但可能把修正分摊到整个动作片段，并降低附近仍然物理可行的动作得分。因此，需要把变化归因到编辑场景中最先发生违规的片段，同时保留其他有效模式。

相关系统分别提供几何感知的有效性、动作流形、动作簇选择、几何感知表示或时空动作条件，但现有记录没有把匹配的最小编辑、首次违规得分归因、边界外锚定和逐模式保留结合起来。Boundary-Local Score Matching (BLSM) 将认证可行性边界处应发生的得分变化，与其他位置不应发生的变化区分开。

近期 VLA 工作和脚本化 RoboTwin 2.0 场景编辑提供了类型化空间定位、持久 3D 上下文、几何引导有效性、可测量的动作模式、成对反事实、模拟器认证的掩码和逐模式得分统计。冻结的连续动作参考策略提供了比较这些得分变化所需的基准。

弥合这一缺口后，VLA 可以在几何编辑后改变特定接触或间隙片段的得分，同时保留其他完成任务的可行方式。这样可以同时通过无效动作抑制、边界外漂移、替代正例覆盖率和最差模式成功率衡量可靠性与多样性。
\section*{方法}
\subsection*{Background}
\emph{构造可接受的几何翻转对，认证其首次违规区间，并枚举受支持的可行模式。}
\begin{enumerate}[leftmargin=*,start=1]
  \item 为每条源轨迹建立结构化记录，字段包括 \(task_{id}\)、任务指令、保存的模拟器状态、冻结策略输出的 T 步连续动作数组及控制频率、原场景重放轨迹和编译器版本。先确认该动作数组在保存状态中成功，再枚举已声明的任务保持型场景编辑；每次都从同一初始状态重放完全相同的动作数组，并按确定性并列规则保留编辑度量 C 最小的结果。【作者需决定：明确可编辑的场景参数、取值域与联动约束，任务身份检查方式，不同编辑在 C 中的归一化方式，以及 C 相同时的选择规则。】用同一个带版本号的模拟器判定规则 F 标记原场景与编辑场景的整段重放。【作者需决定：为每类任务定义成功事件、物理无效事件、数值容差、持续时间要求及其布尔组合。】BoundaryFlip 记录保存两个场景状态、编辑量、C、动作数组、轨迹、工具版本和接纳状态；资格台账则为每次尝试保存判定结果及机器可读的纳入或排除原因。
\noindent\emph{编译器寻找代价最小的几何改动，使同一动作片段从可行变为不可行。}
\par\noindent{\small\ttfamily [eq~1, raw LaTeX, not typeset] \textbackslash\{\}Delta g\textasciicircum{}*=\textbackslash\{\}arg\textbackslash\{\}min\_\{\textbackslash\{\}Delta g\} C(\textbackslash\{\}Delta g)\textbackslash\{\}quad\textbackslash\{\}text\{s.t.\}\textbackslash\{\}quad F(x,a\textasciicircum{}-)=1,\textbackslash\{\};F(x\textbackslash\{\}oplus\textbackslash\{\}Delta g,a\textasciicircum{}-)=0}
\par\emph{为什么：} 固定任务和动作身份后，几何变化就是预期差异，也能显示合格边界对是否超出数据中的某个被挑选角落。
  \item 恢复成对样本共享的初始模拟器状态和随机状态，以相同控制频率在原场景与编辑场景中重放未改动的 T 步动作数组，并在每一步之后记录对齐的状态、接触、约束和任务事件。把编辑场景截至时刻 t 的轨迹前缀映射为二值违规指示 \(v_{t}\)。【作者需决定：定义各任务要监测的事件、接触与状态容差、事件需持续多久、何时解除，以及多个事件如何合成为 \(v_{t}\)。】找到最早违规时刻 \(\tau \)，从该处沿连续违规时刻扩展到第一次解除或数组末尾；仅在这段闭区间内令首段违规掩码 \(m_{t}\) 为一，其余时刻构成补集。若没有任何违规时刻，则在资格台账中把该样本标为 mask-unresolved，不输出空掩码。保存两个长度为 T 的位数组、\(\tau \)、区间终点、支撑事件值和同步轨迹标识。
\noindent\emph{掩码取从最早违规时刻 \(\tau \) 开始、由模拟器确认的首个连续违规区间。}
\par\noindent{\small\ttfamily [eq~2, raw LaTeX, not typeset] \textbackslash\{\}tau=\textbackslash\{\}min\textbackslash\{\}\{t:v\_t(x\textbackslash\{\}oplus\textbackslash\{\}Delta g\textasciicircum{}*,a\textasciicircum{}-\_\{1:t\})=1\textbackslash\{\}\},\textbackslash\{\}qquad m\_t=\textbackslash\{\}mathbf\{1\}\textbackslash\{\}!\textbackslash\{\}left[t\textbackslash\{\}in\textbackslash\{\}mathcal\{R\}(\textbackslash\{\}tau)\textbackslash\{\}right]}
\par\emph{为什么：} 该掩码来自观测到的模拟器结构，确定几何编辑首次改变功能有效性的时间区域。
  \item 从编辑场景保存的初始状态运行任务编译器，使用与无效动作数组相同的任务指令、T 步时域、控制频率和动作表示。枚举编译器声明的解空间，删除完全相同的动作数组，在编辑场景中重放每个候选，只保留通过同一整段判定规则 F 的候选。【作者需决定：明确候选表示、求解器或生成器、声明的解空间、去重规则，以及据以认定枚举完整的穷尽或覆盖证明。】用一个确定性的动作模式定义对保留结果分组。【作者需决定：明确决定模式身份的对象与等价关系，包括所采用的轨迹、接触序列、物体侧别、路径类别或阈值条件。】每个受支持模式只保留一个规范化的待评分动作数组。【作者需决定：定义模式内代表项的选择规则和比较度量，例如采用编译器指定的规范解或已声明度量下的中心样本。】只要某模式至少有一个枚举成员通过 F，就视为受支持。按确定顺序保存替代正例名册 \(A^{+}\)；每条记录包含 \(mode_{id}\)、\(T\times D\) 动作数组、编译器来源、重放结果、代表项依据、模式内成员数和总模式数 K。
\par\emph{为什么：} 按模式编制名册使正例覆盖率和最差模式保留率可见，并防止主导模式聚合隐藏可行替代方案。
\end{enumerate}
\subsection*{BLSM}
\emph{施加边界局部偏好，同时锚定补集和每个可行替代方案。}
\begin{enumerate}[leftmargin=*,start=4]
  \item 将训练前的参考检查点固定在评估模式，并由该检查点初始化可训练副本。在编辑后的观测 \(x\prime \) 下，对无效动作数组和每个模式代表项的每个动作时刻，分别计算可训练策略标量得分 \(s_\theta \) 与参考策略标量得分 \(s_{ref}\)；两者必须使用完全相同的观测预处理和条件上下文。【作者需决定：指定连续动作策略的参数化形式，并定义标量得分的计算方法，包括条件上下文、动作维度归约、归一化和对应的参考计算。】比较前要让每个替代动作的时间索引与首段违规掩码 \(m_{t}\) 所表示的物理时刻一致。【作者需决定：要求所有动作片段共享 T 和控制频率，或定义处理缺失时刻与接触相位偏移的确定性对齐规则。】用可训练得分减去对应参考得分，得到参考校正变化 \(\Delta _\theta \)。对每个模式 k，在 m 内求可行替代动作相对无效动作的平均变化，得到间隔 \(d_\theta ^(k)\)；再使用已规定的平滑逻辑偏好并对 K 个模式等权平均，生成边界损失 \(L_{B}\)。保存 \(L_{B}\)、各 \(mode_{id}\) 及其间隔、掩码时刻数和逐时得分记录。
\noindent\emph{动作得分以相对冻结参考策略的变化量表示，避免原始得分尺度主导偏好。}
\par\noindent{\small\ttfamily [eq~3, raw LaTeX, not typeset] \textbackslash\{\}Delta\_\textbackslash\{\}theta(a\_t\textbackslash\{\}mid x')=s\_\textbackslash\{\}theta(a\_t\textbackslash\{\}mid x')-s\_\{\textbackslash\{\}mathrm\{ref\}\}(a\_t\textbackslash\{\}mid x')}
\noindent\emph{边界损失提升每个可行替代动作相对无效动作的掩码参考校正间隔；\(M_{B}\) 是留出数据上的平均间隔。}
\par\noindent{\small\ttfamily [eq~4, raw LaTeX, not typeset] d\_\textbackslash\{\}theta\textasciicircum{}\{(k)\}=\textbackslash\{\}frac\{\textbackslash\{\}sum\_t m\_t[\textbackslash\{\}Delta\_\textbackslash\{\}theta(a\textasciicircum{}+\_\{k,t\}\textbackslash\{\}mid x')-\textbackslash\{\}Delta\_\textbackslash\{\}theta(a\textasciicircum{}-\_t\textbackslash\{\}mid x')]\}\{\textbackslash\{\}sum\_t m\_t\},\textbackslash\{\}quad \textbackslash\{\}mathcal\{L\}\_B=\textbackslash\{\}frac\{1\}\{K\}\textbackslash\{\}sum\_\{k=1\}\textasciicircum{}\{K\}\textbackslash\{\}log(1+e\textasciicircum{}\{-d\_\textbackslash\{\}theta\textasciicircum{}\{(k)\}\}),\textbackslash\{\}quad M\_B=\textbackslash\{\}mathbb\{E\}[d\_\textbackslash\{\}theta]}
\par\emph{为什么：} 参考调整隔离训练后的移动，而掩码将可行性得分归因指向反事实涉及的片段，而非整个动作片段。
  \item 对无效动作数组和 K 个模式代表项，取首段违规掩码之外的时刻，将每个参考校正得分变化 \(\Delta _\theta \) 平方；先在各动作数组内求平均，再对 K+1 个数组等权平均，得到补集损失 \(L_{comp}\)。若掩码覆盖全部 T 个时刻，该平均没有有效分母。【作者需决定：明确并记录整段掩码样本是被排除、不贡献补集项，还是采用另一种明确定义的处理。】另行对每个可行模式在全部 T 个时刻的 \(\Delta _\theta \) 平方求平均，再对 K 个模式等权平均，得到替代保留损失 \(L_{alt}\)；不得按编译器生成频率或模式内成员数加权。保存两个标量损失，以及可复现各级归约所需的样本级、模式级和逐时统计。
\noindent\emph{补集锚定限制边界外漂移；等权替代项逐一保护所有可行模式，而非只保护主导模式。}
\par\noindent{\small\ttfamily [eq~5, raw LaTeX, not typeset] \textbackslash\{\}mathcal\{L\}\_\{\textbackslash\{\}mathrm\{comp\}\}=\textbackslash\{\}frac\{1\}\{K+1\}\textbackslash\{\}sum\_\{z\textbackslash\{\}in\textbackslash\{\}\{a\textasciicircum{}-,a\_1\textasciicircum{}+,\textbackslash\{\}ldots,a\_K\textasciicircum{}+\textbackslash\{\}\}\}\textbackslash\{\}frac\{\textbackslash\{\}sum\_t(1-m\_t)\textbackslash\{\}Delta\_\textbackslash\{\}theta(z\_t\textbackslash\{\}mid x')\textasciicircum{}2\}\{\textbackslash\{\}sum\_t(1-m\_t)\},\textbackslash\{\}quad \textbackslash\{\}mathcal\{L\}\_\{\textbackslash\{\}mathrm\{alt\}\}=\textbackslash\{\}frac\{1\}\{K\}\textbackslash\{\}sum\_\{k=1\}\textasciicircum{}\{K\}\textbackslash\{\}frac\{1\}\{T\}\textbackslash\{\}sum\_t\textbackslash\{\}Delta\_\textbackslash\{\}theta(a\textasciicircum{}+\_\{k,t\}\textbackslash\{\}mid x')\textasciicircum{}2}
\par\emph{为什么：} 补集项防止时间信用泄漏，替代项使压制少数可行模式变得有代价，即使平均有效性提高。
\end{enumerate}
\subsection*{M3\_validation}
\emph{选择冻结的损失权重，并公开边界、有效性、漂移、覆盖率和最差模式诊断。}
\begin{enumerate}[leftmargin=*,start=6]
  \item 用常规模仿示范和已接纳的 BoundaryFlip 记录构造训练批次，并在常规示范字段上计算基础策略原有的模仿损失。【作者需决定：指定参考检查点、可训练策略类别、动作表示、条件输入、初始化关系、原生模仿损失及批次归约方式。】以非负权重 \(\lambda _{B}\)、\(\lambda _{C}\) 和 \(\lambda _{A}\) 加入边界损失 \(L_{B}\)、补集损失 \(L_{comp}\) 与替代保留损失 \(L_{alt}\)。验证集按原场景/编辑场景成对划分；评估预先声明的权重笛卡尔积，只保留满足补集漂移和最差模式不退化约束的设置，在其中选择边界分离最大的设置，按确定规则处理并列，然后冻结权重。【作者需决定：预先声明三个权重的候选值、两项非劣容差、验证统计量与不确定性判定规则，以及并列处理规则。】留出集边界间隔 \(M_{B}\) 先在每个样本内对模式等权平均，再对样本平均；同时报告尝试数与接纳数及原因、掩码长度及其占 T 的比例、名册大小 K 和留出集补集漂移。通过在编辑场景中执行训练后策略来估计替代正例覆盖率和最差模式保留率。【作者需决定：定义执行所服从的分布、执行结果到编译器模式的映射、模式成功覆盖的标准、保留率采用绝对值还是相对参考策略，以及样本与模式的聚合方式。】保存完整训练配置、数据划分清单、选定权重、模型检查点标识和逐样本诊断记录。输出是用边界局部分数匹配(Boundary-Local Score Matching, BLSM)训练的连续动作视觉---语言---动作(Vision-Language-Action, VLA)策略。
\noindent\emph{BLSM 在常规模仿目标上加入边界偏好、补集稳定与模式保护；三个权重在验证对上一次选定后固定。}
\par\noindent{\small\ttfamily [eq~6, raw LaTeX, not typeset] \textbackslash\{\}mathcal\{L\}=\textbackslash\{\}mathcal\{L\}\_\{\textbackslash\{\}mathrm\{imit\}\}+\textbackslash\{\}lambda\_B\textbackslash\{\}mathcal\{L\}\_B+\textbackslash\{\}lambda\_C\textbackslash\{\}mathcal\{L\}\_\{\textbackslash\{\}mathrm\{comp\}\}+\textbackslash\{\}lambda\_A\textbackslash\{\}mathcal\{L\}\_\{\textbackslash\{\}mathrm\{alt\}\}}
\par\emph{为什么：} 该选择规则避免任意损失权衡冒充机制，而公开的诊断把得分移动与下游可靠性和保留率联系起来。
\end{enumerate}
\end{document}
